% !TEX root = ../projeto-pesquisa.tex
% =====================================================================
%  INTRODUÇÃO
% =====================================================================
\newpage
\secao{Introdução}
\justifying
O avanço das tecnologias de inteligência artificial tem transformado
significativamente a produção e circulação de conteúdos digitais, especialmente
no que se refere à geração de imagens altamente realistas. Esse cenário amplia os
desafios relacionados à verificação da autenticidade de imagens digitais, uma vez
que modelos generativos são capazes de produzir conteúdos visualmente coerentes e
difíceis de distinguir de registros reais. Segundo Padilha \textit{et al.} (2021,
p. 131), o ambiente digital contemporâneo já é marcado por uma elevada presença
de conteúdos manipulados, incluindo montagens, informações fora de contexto e
materiais que dificultam a identificação de sua veracidade.

De acordo com as empresas DeepStrike (2025) e Bright Defense (2026), o volume de
\textit{deepfakes} (imagens, vídeos e áudios sintéticos) compartilhados
\textit{online} explodiu nos últimos anos: passou de aproximadamente 500 mil em
2023 para uma projeção de 8 milhões em 2025, representando um crescimento anual
de cerca de 1.500\%.

Nesse contexto, o presente projeto aborda a análise de padrões e inconsistências
na identificação de imagens geradas por inteligência artificial, com foco em
\textit{deepfakes} e demais formas de manipulação digital. A proposta consiste no
desenvolvimento de uma solução tecnológica voltada à análise de imagens, como uma
plataforma \textit{online}, capaz de apoiar a verificação da autenticidade de
arquivos digitais. Essa necessidade se intensifica diante da crescente
sofisticação dos modelos generativos, que ampliam o potencial de disseminação de
desinformação e golpes digitais no ambiente \textit{online}.

O problema central da pesquisa está relacionado à dificuldade de distinguir
imagens reais de imagens sintéticas, especialmente em casos de manipulação como o
\textit{copy-move}, em que a distribuição estatística dos \textit{pixels} pode
permanecer semelhante à original, dificultando a detecção visual e automatizada.
De acordo com Nagm \textit{et al.} (2024, p. 2), esse tipo de adulteração
representa um desafio relevante na área de forense digital, justamente por
preservar características estatísticas da imagem, o que compromete métodos
tradicionais de detecção baseados apenas em análise visual ou estatística
superficial.

Ainda nesse contexto, técnicas como a Análise de Nível de Erro (ELA) são
utilizadas para evidenciar inconsistências em imagens, por meio da amplificação
de diferenças de compressão que podem indicar possíveis manipulações (Nagm
\textit{et al.}, 2024, p. 4). No entanto, com o avanço dos modelos generativos de
alta fidelidade, a eficácia dessas abordagens tem sido reduzida, exigindo a
integração com métodos mais robustos, como redes neurais convolucionais (CNN),
capazes de identificar padrões complexos de manipulação, além da análise de
metadados EXIF, que fornecem informações técnicas sobre a origem e o histórico de
processamento das imagens (Yang \textit{et al.}, 2026, p. 15).

Estudos recentes mostram que a precisão humana na detecção de \textit{deepfakes}
fica em torno de 50--55\% (próxima ao nível de acaso), e apenas 0,1\% dos
participantes consegue identificar corretamente todos os conteúdos sintéticos e
reais apresentados (iProov, 2025).

Além disso, abordagens contemporâneas indicam que a autenticação de mídias
digitais não deve se basear apenas na análise do conteúdo visual, mas também na
verificação da proveniência e integridade dos arquivos. Segundo England
\textit{et al.} (2020, p. 115), em um cenário próximo, nem humanos nem algoritmos
serão capazes de distinguir de forma confiável conteúdos falsos de conteúdos
reais apenas pela aparência, tornando necessária a adoção de sistemas de
verificação baseados em origem, integridade e rastreabilidade da mídia.

A execução do projeto justifica-se pela necessidade de fortalecer a segurança
digital e promover o pensamento crítico no uso de mídias digitais. Dessa forma, o
projeto contribui para a formação acadêmica dos envolvidos, ao aplicar
conhecimentos de inteligência artificial, ciência de dados e segurança da
informação em um problema real, além de possuir relevância social ao auxiliar na
mitigação da desinformação e na promoção de um ambiente digital mais confiável.

